\scp{Implications from a more complete picture of marsupial terms}{14-06}
Both **kVndoR(a) ‘cuscus’
and **mantəR ‘cuscus’ link in various ways to closely resemblant
words found to the west in Sulawesi
and to the east in Indonesian Papua.
Both **kVndoR(a) ‘cuscus’ and **mantəR ‘cuscus’ show Austronesian
and Papuan cross-over in several places
and with several forms.
Such points of cross-over require that contact
and borrowing are a necessary part of their histories.
Both **kVndoR(a) ‘cuscus’ and **mantəR ‘cuscus’ are in a kind of complementary distribution.
Such complementary distribution in different nodes of the \tsc{Seram-Tanimbar-Bomberai}
subgroup also requires that contact,
rather than inheritance, is a likely explanation of the forms in \tsc{Seram-Tanimbar-Bomberai}.

Consideration of a fuller breadth of the data relating to putative
PCEMP **kandoRa and **mans(aə)r provides little support for the validity of CEMP.
Cognates of **mantəR are found outside the purported “CEMP” area in \tsc{Muna-Buton}
languages, and **kVndoR(a) is only attested unproblematically in EMP
and in a single “CMP” witness (Watubela) on the edge of the area
where EMP languages are spoken.
**kVndoR(a) further appears to be linked to a network of formally
similar forms, such as **kindoR and {\#}kisor(o), cognates of which link incrementally
through Sulawesi, Halmahera, and West Papua, including in non-Austronesian languages.
A form like **kindoR accounts for nearly all the data in the
\tsc{Eastern Islands} node of \tsc{Seram-Tanimbar-Bomberai}
and SHWNG languages, including Watubela.
**kandoRa does not explain most of the related forms in the region.

\citet[262]{Blust2012b} sees his PCEMP
**kandoRa and **mans(aə)r as “a critical piece of evidence [{\ldots}]
for the phylogenetic unity of CEMP”,
although not the sole evidence.
But the full breadth of the on-the-ground evidence for **kVndoR(a)
and **mantəR turn out to actually undermine the CEMP hypothesis.
%rather than being the “jewel in the crown” that establish and support it.